\documentclass[10pt,a4paper]{amsart}
\usepackage{amsmath,amssymb,amsthm,graphicx}
% Retained preamble for the Stacks Project chapter morphisms.tex (Morphisms of Schemes).
% The chapter file pinned-original-source.tex carries only \input{preamble} on its first line, so no
% preamble text lives inside the chapter itself.  The Stacks Project uses one single shared preamble for
% every chapter; the snapshot stages it verbatim in sources/stacks-chow/preamble.tex.  The block below is
% copied mechanically, line for line, from that file's lines 14-18 (the Xy-pic setup) and lines 154-217
% (the native theorem declarations together with the mathematical macros).  Nothing was added, removed or
% redefined.  The wrapper class and the amsmath/amssymb/amsthm/graphicx base packages are supplied by the
% assembler, so the class/\externaldocument part of the shared preamble is not repeated here.
% For commutative diagrams we use Xy-pic
\usepackage[all]{xy}

% We use 2cell for 2-commutative diagrams.
\xyoption{2cell}
% Theorem environments.
%
\theoremstyle{plain}
\newtheorem{theorem}[subsection]{Theorem}
\newtheorem{proposition}[subsection]{Proposition}
\newtheorem{lemma}[subsection]{Lemma}

\theoremstyle{definition}
\newtheorem{definition}[subsection]{Definition}
\newtheorem{example}[subsection]{Example}
\newtheorem{exercise}[subsection]{Exercise}
\newtheorem{situation}[subsection]{Situation}

\theoremstyle{remark}
\newtheorem{remark}[subsection]{Remark}
\newtheorem{remarks}[subsection]{Remarks}

\numberwithin{equation}{subsection}

% Macros
%
\def\lim{\mathop{\mathrm{lim}}\nolimits}
\def\colim{\mathop{\mathrm{colim}}\nolimits}
\def\Spec{\mathop{\mathrm{Spec}}}
\def\Hom{\mathop{\mathrm{Hom}}\nolimits}
\def\Ext{\mathop{\mathrm{Ext}}\nolimits}
\def\SheafHom{\mathop{\mathcal{H}\!\mathit{om}}\nolimits}
\def\SheafExt{\mathop{\mathcal{E}\!\mathit{xt}}\nolimits}
\def\Sch{\mathit{Sch}}
\def\Mor{\mathop{\mathrm{Mor}}\nolimits}
\def\Ob{\mathop{\mathrm{Ob}}\nolimits}
\def\Sh{\mathop{\mathit{Sh}}\nolimits}
\def\NL{\mathop{N\!L}\nolimits}
\def\CH{\mathop{\mathrm{CH}}\nolimits}
\def\proetale{{pro\text{-}\acute{e}tale}}
\def\etale{{\acute{e}tale}}
\def\QCoh{\mathit{QCoh}}
\def\Ker{\mathop{\mathrm{Ker}}}
\def\Im{\mathop{\mathrm{Im}}}
\def\Coker{\mathop{\mathrm{Coker}}}
\def\Coim{\mathop{\mathrm{Coim}}}

% Boxtimes
%
\DeclareMathSymbol{\boxtimes}{\mathbin}{AMSa}{"02}

%
% Macros for moduli stacks/spaces
%
\def\QCohstack{\mathcal{QC}\!\mathit{oh}}
\def\Cohstack{\mathcal{C}\!\mathit{oh}}
\def\Spacesstack{\mathcal{S}\!\mathit{paces}}
\def\Quotfunctor{\mathrm{Quot}}
\def\Hilbfunctor{\mathrm{Hilb}}
\def\Curvesstack{\mathcal{C}\!\mathit{urves}}
\def\Polarizedstack{\mathcal{P}\!\mathit{olarized}}
\def\Complexesstack{\mathcal{C}\!\mathit{omplexes}}
% \Pic is the operator that assigns to X its picard group, usage \Pic(X)
% \Picardstack_{X/B} denotes the Picard stack of X over B
% \Picardfunctor_{X/B} denotes the Picard functor of X over B
\def\Pic{\mathop{\mathrm{Pic}}\nolimits}
\def\Picardstack{\mathcal{P}\!\mathit{ic}}
\def\Picardfunctor{\mathrm{Pic}}
\def\Deformationcategory{\mathcal{D}\!\mathit{ef}}

% --- appended 2026-10-12 (Chief): the project own definitions of the reference/slogan/history
% display environments, copied mechanically from the shared Stacks preamble lines 8-12
% (\usepackage{verbatim} plus the three \newenvironment lines, which the chapter preamble keeps
% verbatim from the project preamble). Nothing added, removed or redefined: these environments
% carry no mathematical content and the project itself typesets them as comments. Without these
% three lines a counted statement whose own text uses them fails with Environment slogan undefined.
\usepackage{verbatim}
\newenvironment{reference}{\comment}{\endcomment}
\newenvironment{slogan}{\comment}{\endcomment}
\newenvironment{history}{\comment}{\endcomment}
\usepackage{cleveref}
\crefname{originallemma}{original lemma}{original lemmas}
\Crefname{originallemma}{Original lemma}{Original lemmas}
\makeatletter
\newlabel{algebra-lemma-normal-ring-integrally-closed}{{034M}{source}{}{sourcepointer-1}{}}
\newlabel{algebra-lemma-normal-ring-integrally-closed@cref}{{[originallemma][0][]034M}{[1][0][]source}{}{}{}}
\makeatother
\crefname{originalequation}{original equation}{original equations}
\Crefname{originalequation}{Original equation}{Original equations}
\makeatletter
\newlabel{equation-generic-points}{{035M}{source}{}{sourcepointer-2}{}}
\newlabel{equation-generic-points@cref}{{[originalequation][0][]035M}{[1][0][]source}{}{}{}}
\makeatother
\crefname{originallemma}{original lemma}{original lemmas}
\Crefname{originallemma}{Original lemma}{Original lemmas}
\makeatletter
\newlabel{lemma-characterize-normalization}{{035I}{source}{}{sourcepointer-3}{}}
\newlabel{lemma-characterize-normalization@cref}{{[originallemma][0][]035I}{[1][0][]source}{}{}{}}
\makeatother
\crefname{originallemma}{original lemma}{original lemmas}
\Crefname{originallemma}{Original lemma}{Original lemmas}
\makeatletter
\newlabel{lemma-description-normalization}{{035P}{source}{}{sourcepointer-4}{}}
\newlabel{lemma-description-normalization@cref}{{[originallemma][0][]035P}{[1][0][]source}{}{}{}}
\makeatother
\crefname{originallemma}{original lemma}{original lemmas}
\Crefname{originallemma}{Original lemma}{Original lemmas}
\makeatletter
\newlabel{lemma-integral-universally-closed}{{01WM}{source}{}{sourcepointer-5}{}}
\newlabel{lemma-integral-universally-closed@cref}{{[originallemma][0][]01WM}{[1][0][]source}{}{}{}}
\makeatother
\crefname{originallemma}{original lemma}{original lemmas}
\Crefname{originallemma}{Original lemma}{Original lemmas}
\makeatletter
\newlabel{lemma-normal-normalization}{{035L}{source}{}{sourcepointer-6}{}}
\newlabel{lemma-normal-normalization@cref}{{[originallemma][0][]035L}{[1][0][]source}{}{}{}}
\makeatother
\crefname{originallemma}{original lemma}{original lemmas}
\Crefname{originallemma}{Original lemma}{Original lemmas}
\makeatletter
\newlabel{lemma-normalization-reduced}{{035O}{source}{}{sourcepointer-7}{}}
\newlabel{lemma-normalization-reduced@cref}{{[originallemma][0][]035O}{[1][0][]source}{}{}{}}
\makeatother
\crefname{originallemma}{original lemma}{original lemmas}
\Crefname{originallemma}{Original lemma}{Original lemmas}
\makeatletter
\newlabel{schemes-lemma-map-into-reduction}{{0356}{source}{}{sourcepointer-8}{}}
\newlabel{schemes-lemma-map-into-reduction@cref}{{[originallemma][0][]0356}{[1][0][]source}{}{}{}}
\makeatother
\title{Structure of the normalization of a base in a normal source with locally finitely many components}
\author{The Stacks Project Authors}
\date{stacks-project:morphisms.tex:local-snapshot-20261012; GNU FDL 1.2 or later}
\begin{document}
\maketitle
\noindent\textit{Editorial note.} All delivered mathematical text is the author's own, copied from the pinned Stacks Project chapter file pinned-original-source.tex; no statement or proof step was rewritten, reordered or removed.\par\medskip
\noindent\textit{Editorial note.} Following the multi-claim convention for this source, source\_url is the source-level URL https://stacks.math.columbia.edu/ and the source label lemma-normalization-normal is carried in claim\_key and here rather than in a per-claim tag URL.\par\medskip
\noindent\textit{Editorial note.} The preamble is the chapter's shared Stacks Project preamble, retained as sources/stacks-morphisms/retained-preamble.tex: a line-for-line excerpt of the shared Stacks preamble staged with this snapshot. No macro was invented or redefined, and this card adds no preamble text of its own.\par\medskip
\noindent\textit{Editorial note.} Statement span 14503-14525 and proof span 14527-14585 come from the halfA pool card operations/pool-cards/stacks-morphisms-halfA.json and were re-read in the pinned chapter file before this card was written.\par\medskip
\noindent\textit{Editorial note.} The proof ends with the author's sentence 'We omit the verification that it is unique.' for the uniqueness part; that sentence is kept verbatim, and the construction of the canonical factorization is carried out in the counted span.\par\medskip
\noindent\textit{Editorial note.} This card was previously skipped because the shared gate substring-matched my chapter name against the published items of more-morphisms.tex; after the Lead repaired the gate the candidate was re-gated (RESULT=OK) and opened.\par\medskip
\noindent\textit{Editorial note.} 8 cross-references occur inside the counted spans; they are listed in external\_references as navigation pointers only.\par\medskip
\noindent\textit{Allowed standard prerequisites (editorial).}\begin{itemize}
\item The structure statement for the normalization in a normal source, established earlier in the same chapter (and established by a previous card of this lane). Cited at proof lines 14527-14531 and 14537-14539. Gives that the normalization is a disjoint union of normal integral schemes and that the intermediate morphism induces a bijection on irreducible components.
\item The characterization of the normalization morphism as integral, and the compatibility of normalization with passing to the reduction, both cited from the same chapter. Cited at proof lines 14536-14537 and 14565-14566. Give integrality of the normalization and let the uniqueness statement be reduced to the reduced case.
\item An integral morphism is universally closed, and the description of the normalization by integral closures and by the construction through a dominant morphism from a normal source, cited from the same chapter. Cited at proof lines 14542-14544, 14553-14555 and 14570-14570. Turn dominance into surjectivity, give the factorization statement, and supply the local algebra used for the uniqueness statement.
\item A morphism to a reduced scheme factors through the reduction, and a normal ring equals its integral closure in its total ring of fractions. Cited at proof lines 14566-14568 and 14580-14582. Reduce to a reduced base and identify the target of the integral closure map with the normal coordinate ring.
\end{itemize}\medskip
\ifdefined\hypertarget\hypertarget{sourcepointer-1}{}\fi
\noindent\textit{Source reference.} Original lemma 034M: focused-seeds/AGP/sources/stacks-snapshot-20261011/tags.txt, source lines 4097--4097. Same-chapter reference used in the retained author text at source line(s) 14579. It is supplied only as a navigation pointer so that the retained text compiles; the counted item delivers the author's own statement and the author's complete proof, and no part of the referenced item is copied in or substituted for the counted argument.\par\medskip
\ifdefined\hypertarget\hypertarget{sourcepointer-2}{}\fi
\noindent\textit{Source reference.} Original equation 035M: focused-seeds/AGP/sources/stacks-snapshot-20261011/tags.txt, source lines 4133--4133. Same-chapter reference used in the retained author text at source line(s) 14387, 14406, 14413, 14462, 14528, 14551, 14996. It is supplied only as a navigation pointer so that the retained text compiles; the counted item delivers the author's own statement and the author's complete proof, and no part of the referenced item is copied in or substituted for the counted argument.\par\medskip
\ifdefined\hypertarget\hypertarget{sourcepointer-3}{}\fi
\noindent\textit{Source reference.} Original lemma 035I: focused-seeds/AGP/sources/stacks-snapshot-20261011/tags.txt, source lines 4129--4129. Same-chapter reference used in the retained author text at source line(s) 14027, 14033, 14059, 14409, 14538, 14556, 14913, 15004. It is supplied only as a navigation pointer so that the retained text compiles; the counted item delivers the author's own statement and the author's complete proof, and no part of the referenced item is copied in or substituted for the counted argument.\par\medskip
\ifdefined\hypertarget\hypertarget{sourcepointer-4}{}\fi
\noindent\textit{Source reference.} Original lemma 035P: focused-seeds/AGP/sources/stacks-snapshot-20261011/tags.txt, source lines 4136--4136. Same-chapter reference used in the retained author text at source line(s) 14491, 14535, 14570, 14599, 14678, 14700, 15083. It is supplied only as a navigation pointer so that the retained text compiles; the counted item delivers the author's own statement and the author's complete proof, and no part of the referenced item is copied in or substituted for the counted argument.\par\medskip
\ifdefined\hypertarget\hypertarget{sourcepointer-5}{}\fi
\noindent\textit{Source reference.} Original lemma 01WM: focused-seeds/AGP/sources/stacks-snapshot-20261011/tags.txt, source lines 2487--2487. Same-chapter reference used in the retained author text at source line(s) 10976, 11172, 11375, 11380, 11512, 11552, 11624, 11672, 11769, 14152, 14165, 14544. It is supplied only as a navigation pointer so that the retained text compiles; the counted item delivers the author's own statement and the author's complete proof, and no part of the referenced item is copied in or substituted for the counted argument.\par\medskip
\ifdefined\hypertarget\hypertarget{sourcepointer-6}{}\fi
\noindent\textit{Source reference.} Original lemma 035L: focused-seeds/AGP/sources/stacks-snapshot-20261011/tags.txt, source lines 4132--4132. Same-chapter reference used in the retained author text at source line(s) 14532, 14539. It is supplied only as a navigation pointer so that the retained text compiles; the counted item delivers the author's own statement and the author's complete proof, and no part of the referenced item is copied in or substituted for the counted argument.\par\medskip
\ifdefined\hypertarget\hypertarget{sourcepointer-7}{}\fi
\noindent\textit{Source reference.} Original lemma 035O: focused-seeds/AGP/sources/stacks-snapshot-20261011/tags.txt, source lines 4135--4135. Same-chapter reference used in the retained author text at source line(s) 14534, 14561, 14597. It is supplied only as a navigation pointer so that the retained text compiles; the counted item delivers the author's own statement and the author's complete proof, and no part of the referenced item is copied in or substituted for the counted argument.\par\medskip
\ifdefined\hypertarget\hypertarget{sourcepointer-8}{}\fi
\noindent\textit{Source reference.} Original lemma 0356: focused-seeds/AGP/sources/stacks-snapshot-20261011/tags.txt, source lines 4117--4117. Same-chapter reference used in the retained author text at source line(s) 291, 987, 10250, 13640, 14408, 14563. It is supplied only as a navigation pointer so that the retained text compiles; the counted item delivers the author's own statement and the author's complete proof, and no part of the referenced item is copied in or substituted for the counted argument.\par\medskip
\begin{lemma}
\label{lemma-normalization-normal}
Let $X$ be a scheme such that every quasi-compact open has
finitely many irreducible components.
\begin{enumerate}
\item The normalization $X^\nu$ is a disjoint union of integral normal schemes.
\item The morphism $\nu : X^\nu \to X$ is integral, surjective, and
induces a bijection on irreducible components.
\item For any integral morphism $\alpha : X' \to X$ such that for
$U \subset X$ quasi-compact open the inverse image $\alpha^{-1}(U)$ has
finitely many irreducible components and
$\alpha|_{\alpha^{-1}(U)} : \alpha^{-1}(U) \to U$ is birational\footnote{This
awkward formulation is necessary as we've only defined what
it means for a morphism to be birational if the source and target
have finitely many irreducible components. It suffices if
$X'_{red} \to X_{red}$ satisfies the condition.}
there exists a factorization
$X^\nu \to X' \to X$ and $X^\nu \to X'$ is the normalization of $X'$.
\item For any morphism $Z \to X$ with $Z$ a normal scheme
such that each irreducible component of $Z$ dominates an irreducible
component of $X$ there exists a unique factorization $Z \to X^\nu \to X$.
\end{enumerate}
\end{lemma}

\begin{proof}
Let $f : Y \to X$ be as in (\ref{equation-generic-points}).
The scheme $X^\nu$ is a disjoint union of normal integral schemes
because $Y$ is normal and every affine open of $Y$ has finitely
many irreducible components, see
Lemma \ref{lemma-normal-normalization}. This proves (1).
Alternatively one can deduce (1) from
Lemmas \ref{lemma-normalization-reduced} and
\ref{lemma-description-normalization}.

\medskip\noindent
The morphism $\nu$ is integral by Lemma \ref{lemma-characterize-normalization}.
By Lemma \ref{lemma-normal-normalization} the
morphism $Y \to X^\nu$ induces a bijection on irreducible components,
and by construction of $Y$ this implies that $X^\nu \to X$ induces
a bijection on irreducible components. By construction $f : Y \to X$
is dominant, hence also $\nu$ is dominant. Since an integral morphism is
closed (Lemma \ref{lemma-integral-universally-closed}) this implies that
$\nu$ is surjective. This proves (2).

\medskip\noindent
Suppose that $\alpha : X' \to X$ is as in (3). It is clear that
$X'$ satisfies the assumptions under which the normalization
is defined. Let $f' : Y' \to X'$ be the morphism
(\ref{equation-generic-points}) constructed starting with $X'$.
As $\alpha$ is locally birational it is clear that
$Y' = Y$ and $f = \alpha \circ f'$.
Hence the factorization $X^\nu \to X' \to X$ exists
and $X^\nu \to X'$ is the normalization of $X'$ by
Lemma \ref{lemma-characterize-normalization}. This proves (3).

\medskip\noindent
Let $g : Z \to X$ be a morphism whose domain is a normal scheme
and such that every irreducible component dominates an irreducible
component of $X$. By Lemma \ref{lemma-normalization-reduced}
we have $X^\nu = X_{red}^\nu$ and by
Schemes, Lemma \ref{schemes-lemma-map-into-reduction}
$Z \to X$ factors through $X_{red}$. Hence we may replace $X$ by
$X_{red}$ and assume $X$ is reduced. Moreover, as the factorization
is unique it suffices to construct it locally on $Z$.
Let $W \subset Z$ and $U \subset X$ be affine opens
such that $g(W) \subset U$. Write $U = \Spec(A)$ and
$W = \Spec(B)$, with $g|_W$ given by $\varphi : A \to B$.
We will use the results of Lemma \ref{lemma-description-normalization} freely.
Let $\mathfrak p_1, \ldots, \mathfrak p_t$ be the minimal primes of $A$.
As $Z$ is normal, we see that $B$ is a normal
ring, in particular reduced. Moreover, by assumption any minimal
prime $\mathfrak q \subset B$ we have that $\varphi^{-1}(\mathfrak q)$
is a minimal prime of $A$. Hence if $x \in A$ is a nonzerodivisor, i.e.,
$x \not \in \bigcup \mathfrak p_i$, then $\varphi(x)$ is a nonzerodivisor
in $B$. Thus we obtain a canonical ring map $Q(A) \to Q(B)$. As $B$ is
normal it is equal to its integral closure in $Q(B)$ (see
Algebra, Lemma \ref{algebra-lemma-normal-ring-integrally-closed}).
Hence we see that the integral closure $A' \subset Q(A)$ of $A$
maps into $B$ via the canonical map $Q(A) \to Q(B)$.
Since $\nu^{-1}(U) = \Spec(A')$ this gives the canonical
factorization $W \to \nu^{-1}(U) \to U$ of $\nu|_W$.
We omit the verification that it is unique.
\end{proof}

\end{document}
